\section*{الفصل \ref{ch:b2:measurement-statistics} --- إحصاء القياس الكيميائي}

\begin{solution}{exo:b2:measurement-statistics:1}
$\bar V = 62.41/5 = \qty{12.482}{mL}$؛ $\sum(V_i - \bar V)^2 = 0.00328$، $s = \sqrt{0.00328/4} \approx
\qty{0.029}{mL}$؛ $s/\sqrt5 \approx \qty{0.013}{mL}$.
\end{solution}

\begin{solution}{exo:b2:measurement-statistics:2}
$t = 2.776$ من أجل 4 درجات حرية: $12.482 \pm 2.776 \times 0.0128$، أي $(12.48 \pm 0.04)$~mL.
\end{solution}

\begin{solution}{exo:b2:measurement-statistics:3}
الارتيابان النسبيان للكتلة وللحجم 0.04~\% و0.06~\%؛ وبجمعهما
تربيعيًا،
\[\frac{u(c)}{c} = \sqrt{\Bigl(\frac{0.0002}{0.5000}\Bigr)^2 + \Bigl(\frac{0.15}{250.0}\Bigr)^2}
\approx 0.07~\%.\]
\end{solution}

\begin{solution}{exo:b2:measurement-statistics:4}
$\bar x = 1.5$، $\bar y = 0.31$، $S_{xx} = 5$، $S_{xy} = 0.99$: $b = 0.198$، $a = 0.31 - 0.198 \times 1.5
= 0.013$.
\end{solution}

\begin{solution}{exo:b2:measurement-statistics:5}
$|10.12 - 10.00|/(0.08/\sqrt5) = 0.12/0.036 \approx 3.4 > 2.776$: الفرق معنوي، 
والطريقة منحازة (بنحو \qty{+0.12}{mg/L}).
\end{solution}

\begin{solution}{exo:b2:measurement-statistics:6}
$\mathrm{pH} = -\log_{10}c$، $\mathrm d\,\mathrm{pH}/\mathrm dc = -1/(c\ln10)$:
$u(\mathrm{pH}) = (u(c)/c)/\ln10 = 0.02/2.303 \approx 0.009$.
\end{solution}

\begin{solution}{exo:b2:measurement-statistics:7}
$x_{\mathrm{LOD}} = 3 \times 0.0012/0.0098 \approx \qty{0.37}{\micro g/L}$؛ $x_{\mathrm{LOQ}} = 10
\times 0.0012/0.0098 \approx \qty{1.2}{\micro g/L}$.
\end{solution}

\begin{solution}{exo:b2:measurement-statistics:8}
يعطي المستقيم \qty{3.11}{mg/L} في المحاليل المقيسة؛ وكانت العينة قد خُففت خمس مرات:
$3.11 \times 50.0/10.0 \approx \qty{15.6}{mg/L}$.
\end{solution}

\begin{solution}{exo:b2:measurement-statistics:9}
القيمة \qty{0.05}{mg/L} هي \omterm{def:b2:measurement-statistics:validation}{قابلية التكرار}، و\qty{0.15}{mg/L} \omterm{def:b2:measurement-statistics:validation}{قابلية الاستنساخ}. فلكل مختبر
انحيازه الصغير الخاص (المعايرة، والمعايير، والجهاز)؛ وتختلف هذه الانحيازات بين المختبرات وتضاف
إلى التشتت، فتكون \omterm{def:b2:measurement-statistics:validation}{قابلية الاستنساخ} دائمًا الأكبر.
\end{solution}

\begin{solution}{exo:b2:measurement-statistics:10}
$\partial S/\partial a = 0$ يُقرأ $\sum(y_i - a - bx_i) = 0$: فمجموع \omterm{def:b2:measurement-statistics:calibration}{البواقي} معدوم. وبالقسمة على
$n$: $\bar y = a + b\bar x$، إذن $(\bar x, \bar y)$ على المستقيم.
\end{solution}

\begin{solution}{exo:b2:measurement-statistics:11}
خطوة واحدة:
\[\sqrt{(0.006/1.000)^2 + (0.08/100.0)^2} \approx 0.61~\%.\]
خطوة عشرية واحدة:
\[\sqrt{(0.02/10.00)^2 + (0.08/100.0)^2} \approx 0.22~\%,\]
وخطوتان منها مجموعتان تربيعيًا تعطيان $0.22\sqrt2 \approx 0.30~\%$. فطريق الخطوتين أدق بمرتين: والماصة الصغيرة هي الغالبة في
طريق الخطوة الواحدة.
\end{solution}

\begin{solution}{exo:b2:measurement-statistics:12}
$z = 1.6/\sqrt{0.4^2 + 0.5^2} \approx 2.5 > 2$: غير متوافقتين؛ فلمختبر واحد (على الأقل) \omterm{def:b2:measurement-statistics:validation}{انحياز}.
ومع الارتيابات الموسعة ($k = 2$)، يعطي الأول $9.6 \pm 0.8$، الذي يشمل 10؛ والثاني
$11.2 \pm 1.0$، الواقع كله فوق 10. وقبل أي حكم على الماء، يجب حل هذا الخلاف،
بمادة مرجعية مثلًا.
% ledger: who:lead
\end{solution}

\begin{solution}{pb:b2:measurement-statistics:1}
\textbf{1.} الحمض يغيّر التذرير والاستجابة؛ فيجب أن تكون للمعايير والعينات المصفوفة
نفسها كي ينطبق الميل.
\textbf{2.} $\bar x = 10$، $\bar y = 0.1004$، $S_{xx} = 250$، $S_{xy} = 2.445$: $b = 0.00978$ لكل
\unit{\micro g/L}، $a = 0.1004 - 0.0978 = 0.0026$.
\textbf{3.} $+0.0004$، $-0.0005$، $+0.0006$، $-0.0013$، $+0.0008$.
\textbf{4.} الإشارات متناوبة، ولا انحناء: فالمستقيم كافٍ على المجال 0--\qty{20}{\micro g/L}.
\textbf{5.} $s_{y/x} = \sqrt{3.1 \times 10^{-6}/3} \approx 0.00102$.
\textbf{6.} الشاهد (الحمض، والماء، والفرن) يعطي امتصاصية صغيرة خاصة به.
\textbf{7.} الحساسية: الامتصاصية لكل \unit{\micro g/L} من الرصاص.
\textbf{8.} $\bar y_0 = 0.1010$؛ $s = 0.0030$.
\textbf{9.} $x_0 = (0.1010 - 0.0026)/0.00978 \approx \qty{10.06}{\micro g/L}$.
\textbf{10.} $s_{x_0} = (0.00102/0.00978)\sqrt{1/3 + 1/5 + (0.1010 - 0.1004)^2/(0.00978^2 \times 250)}
\approx 0.104 \times 0.730 \approx \qty{0.076}{\micro g/L}$.
\textbf{11.} $n - 2 = 3$؛ $t = 3.182$.
\textbf{12.} $10.06 \pm 3.182 \times 0.076 = (10.06 \pm 0.24)$~\unit{\micro g/L}.
\textbf{13.} ينخفض الحد $1/m$ من 1 إلى $1/3$: فينخفض $s_{x_0}$ بنحو الثلث.
\textbf{14.} $f = 50.50/50.0 = 1.010$؛ $10.06 \times 1.010 \approx \qty{10.16}{\micro g/L}$.
\textbf{15.} $f = 1 + V_2/V_1$: $u^2(f) = (u(V_2)/V_1)^2 + (V_2u(V_1)/V_1^2)^2 = (0.005/50.0)^2 +
(0.50 \times 0.03/2500)^2$، $u(f) \approx 1.0 \times 10^{-4}$، أي 0.01~\% نسبيًا.
\textbf{16.} لا: 0.01~\% مقابل $0.076/10.06 \approx 0.75~\%$.
\textbf{17.} $s_{\mathrm{blank}} \approx 0.00115$؛ $x_{\mathrm{LOD}} = 3 \times 0.00115/0.00978 \approx
\qty{0.35}{\micro g/L}$.
\textbf{18.} $10 \times 0.00115/0.00978 \approx \qty{1.2}{\micro g/L}$.
\textbf{19.} نعم، بفارق كبير.
\textbf{20.} نعم: من 9.92 إلى \qty{10.40}{\micro g/L} في ماء الصنبور.
\textbf{21.} $(0.104 - 0.0026)/0.00978 \times 1.010 \approx \qty{10.5}{\micro g/L}$، وكانت ستعطي، وحدها،
حكمًا واثقًا بأنها «فوق القيمة الإرشادية» لا تدعمه المعطيات.
\textbf{22.} القاعدة المستعملة لبناء المجال تلتقط القيمة الحقيقية في 95~\% من السلاسل التي
تُطبق عليها؛ وليست احتمالًا يخص هذه القيمة الواحدة.
\textbf{23.} مزيد من القراءات المكررة ومزيد من المعايير قرب \qty{10}{\micro g/L} (الحدان $1/m$ 
و$1/n$)، أو طريقة أكثر حساسية؛ والمجال لا يضيق إلا ببطء، كالمقدار $1/\sqrt{\text{العدد}}$.
\textbf{24.} بتحليل ماء مرجعي معتمد، أو بإضافة كمية معلومة من الرصاص إلى العينة
والتحقق من نسبة الاسترداد.
\textbf{25.} قرب \qty{10}{\micro g/L} يكون ارتياب الطرق الروتينية بضعة بالمئة، كما هنا: فقيمة إرشادية
أدنى ما كان يمكن التحقق منها على نحو موثوق في المختبرات العادية.
\textbf{26.} $\boldsymbol{(10.16 \pm 0.24)}$~\unit{\micro g/L} عند 95~\%: المجال يحوي
\qty{10}{\micro g/L}، فلا يبيّن القياس أن الماء يتجاوز القيمة الإرشادية ولا أنه يستوفيها
بهامش.
\end{solution}
